\section{Introduction and scope of the results}\label{sec:introduction}

An arrangement of $n$ straight lines divides the real plane into bounded and unbounded cells. The Kobon problem asks how many of the bounded cells can be triangles. The difficulty is not creating many triples of intersecting lines: there are $\binom n3$ candidate triples. The difficulty is ensuring that their triangles are not cut by any of the remaining lines and that their interiors do not overlap. The distinction between an attractive picture and a valid arrangement is therefore central to both construction and verification.

Let $\K(n)$ denote the classical maximum, allowing multiple intersections and parallel lines, and let $\Ks(n)$ denote the maximum over simple arrangements: no parallel pair and no three concurrent lines. All lower-bound witnesses constructed by the universal theorem below are simple, so they establish lower bounds for both quantities. Some of the strongest retained finite certificates are nonsimple. We never use an upper bound restricted to $\Ks$ as an upper bound for $\K$ without an additional argument.

This paper consolidates the author's March 2026 extension proposal and the subsequent program of exact construction, comparison with the literature, and formal proof. The purpose is both mathematical and documentary: to state the strongest completed results, to make their proofs and provenance inspectable, and to preserve the failed approaches that determine the next research questions. The original general extension claim is reported as an unresolved claim rather than silently promoted to a theorem.

\subsection{The main completed all-order theorem}

Define $\B(n)=\ceil{n(n-3)/3}$ for $n\ge3$, and set $\B(n)=0$ for $n<3$. The familiar F\"uredi--Pal\'asti construction supplies this affine baseline; the underlying construction and its historical priority belong to the earlier literature~\cite{furedipalasti1984,erickson1999}. The first completed all-order theorem is the following refinement; the recursive envelope below now strengthens its retained finite enhancement.

\begin{theorem}[Universal affine refinement]\label{thm:main}
For every $n\in\N$, there is a simple arrangement of $n$ real straight lines with at least $\G(n)$ bounded triangular cells, where
\begin{equation}\label{eq:G}
\G(n)=\begin{cases}
0,&0\le n<3,\\
1,&n=3,\\
\floor{n(n-3)/3}+1+(n\bmod2),&n\ge4.
\end{cases}
\end{equation}
Consequently $\K(n)\ge\Ks(n)\ge\G(n)$.
\end{theorem}

The proof, developed in Section~\ref{sec:universal}, combines a shifted trigonometric arrangement with one- and two-cap changes of affine chart. The Lean declaration \lean{Kobon.Universal.baseline_sound} proves the geometric existence statement for arbitrary $n$. It has no assumed extension theorem and no dependence on finite experimental verification. Its dependencies are the standard logical axioms used by Lean and Mathlib.

For $n\ge4$, the gains over $\B(n)$ are $1,1,0,2,0,1$ in residue classes $0,1,2,3,4,5$ modulo six. Thus the refinement is strict at every odd order at least five and every positive multiple of six. Both formulas have leading behavior $n^2/3-n+O(1)$: this is a constant-term refinement, not an improved quadratic or linear asymptotic coefficient. The original F\"uredi--Pal\'asti projective count table already contains related constants. The present source audit does not establish first numerical priority for the affine formula.

\begin{figure}[tbp]
\centering\includegraphics[width=0.94\linewidth]{bounds-growth.pdf}
\caption{Growth of the classical general baseline, the proved affine refinement, and the classical polynomial upper benchmark. The close overlap on a quadratic scale is mathematically meaningful: the new uniform gain is bounded. Residue and deficit plots below resolve the smaller differences. This figure is computed from the formulas, not copied from earlier publications.}\label{fig:growth}
\end{figure}

\subsection{Finite enhancements and other completed components}

The universal construction is supplemented by exact coordinates. Let $C(n)$ be the largest triangle count among the finite certificates in the release at order $n$, with $C(n)=0$ when no such certificate is saved. Then
\begin{equation}\label{eq:envelope}
\Lbound(n)=\max\{\G(n),C(n)\},\qquad \K(n)\ge\Lbound(n)
\end{equation}
is proved for all natural orders. This elementary maximum operation retains published constructions and the author's earlier values without altering their attribution. It is not a claim that $\Lbound$ is the largest lower bound obtainable from all published constructions by every possible interpolation or extension.

The finite catalog contains 138 distinct coordinate certificates, including 104 simple witnesses. A separate generated table has 64 orders at which a saved finite certificate strictly improves the older baseline $\B$; some of these improvements are already supplied by $\G$. The catalog includes weaker checkpoints and independent reproductions. Its size is not a count of original discoveries.

The 2 October release added a recursively closed geometric family. Put
$q_t=10\cdot2^t$ and $A_t=(q_t^2-4)/3$. For every $t\ge0$, simple
arrangements realize at least $A_t$ triangles at $q_t+1$ lines and
$A_t+q_t/2$ at $q_t+2$ lines. Every next seed retains the sorted grid,
saturation, positive central apex and visible-boundary data. Taking
the maximum of these candidates and $\Lbound$ gives the earlier total
bound $H$ of Theorem~\ref{fam:all-n-envelope}. It strictly improves
the previous repository envelope at every family pair from 321/322
onward, without claiming a new literature record for the inherited
numerical family.

The research of 5--6 October completes the compatible 33- and 49-line
seed proofs, adds the known 37-line orbit, and realizes a further
60-dyadic orbit from Poola's 61-line point input. Its strongest even
branch uses 29 visible pairs. The resulting portfolio $J$ retains $H$
pointwise (Theorem~\ref{fam:portfolio}). These are actual straight-line
witnesses at every finite depth; their historical attributions and
simple-model gaps appear in Table~\ref{fam:orbit-table}.

The earlier manuscript's quantitative boundary-defect successor now has
an end-to-end Lean proof, including terminal-ray extraction and averaging,
and it iterates through both parities. The original full half-order
recurrence remains unproved. On the upper side, actual cyclic fans,
side matching, graph components and geometric escape are now extracted
from the arrangement. The strongest all-triple global estimate is
\[
 2\delta+B_{15}\ge2U+3(E_3+P_{21})+N_{12}+M_{22},
\]
with the local counts defined in Section~\ref{upper:degree-free}.
Mixed-multiplicity and component estimates retain explicit corrections;
they do not yield an unrestricted numerical bound depending only on $n$.
The translation analysis of Section~\ref{sec:surgery} proves exact
birth/loss accounting and counterexamples to a specified loss formula.
Table~\ref{tab:claims} separates these claims from the still open problem.

\begin{table}[tbp]
\centering\small
\begin{tabularx}{\linewidth}{@{}P{0.28\linewidth}XP{0.21\linewidth}@{}}
\toprule
Result & Mathematical content & Release status\\
\midrule
Universal $\G$ & Simple real-line witnesses at every order & General Lean proof\\
Envelope $\Lbound$ & Maximum with all saved finite witnesses & Lean; finite native checks\\
28:238, 30:275, 34:357 & Earlier attributed finite inequalities & Exact Lean certificates\\
Exterior extension & New line and gain from a visible-pair list & General Lean proof\\
Full successor recurrence & $\K(n+1)\ge\K(n)+\floor{n/2}$ & Unproved\\
49-seed numerical target & Lower function obtained by summing the proposed recurrence & Independently proved\\
One-step BBL doubling & $(q+1,T)\mapsto(2q+1,T+q^2)$ for compatible seeds & General Lean proof\\
Infinite hybrid iteration & Odd and full-gain even geometric families & Complete Lean; native seed\\
33/37/49 orbits & Known numerical families, compatible real seeds and simple optimality & Complete Lean; native seeds\\
61 orbit & Uniform $1190/29$ seed; constant simple gaps 9 and 10 & Complete Lean; native seed\\
Quantitative successor & Deficit-dependent actual rule, both parities, indefinite iteration & General Lean proof\\
Local fan structure & Incidence, antipodal and run-length restrictions & General Lean proofs\\
Simple upper windows & Actual odd/even upper bounds and family existence & Complete Lean\\
Actual clean-line charging & Charge existence and finite fiber bounds & General Lean proof\\
Global multiplicity consequences & Actual fans, matching, mixed and triple defect estimates & General Lean proofs; class scope\\
Local surgery & Exact translation germs, retention, births and losses & General Lean proofs\\
Loss-rule counterexamples & Four- and six-line real parameter intervals & General Lean proofs\\
Tangent-grid obstruction & Explicit ten-sign algebraic impossibility & General Lean proof\\
\bottomrule
\end{tabularx}
\caption{Claim status in the 7 October edition. ``General Lean proof'' does not imply that every proposed application has been formalized. Published construction inputs retain their original authorship.}\label{tab:claims}
\end{table}

\subsection{History, attribution, and the structure of the paper}

The March archive records a proposed extension from odd to even orders, associated finite bounds, and an incomplete Lean file. That file contains eight proof placeholders, including its central extension step. Later exact certificates establish the finite inequalities independently, while the current all-order proof uses a different construction. Accordingly, a corrected account must separate the survival of numerical conclusions from the validity of the original proof strategy.

The September work broadened the program to both parities, exterior visibility, boundary-defect estimates, compatible doubling seeds, exact rational realization, and multiplicity-sensitive upper arguments. Several plausible routes failed: direct repetition of exterior gains lost the necessary boundary resources; small perturbations of nonsimple records sometimes reduced the triangle count; and large local searches failed to improve their chosen seeds. These negative findings retain their finite scope. The 2 October work closed the recursive and clean-line interfaces and added grid obstructions. The 5--6 October work closes the quantitative successor and actual nonsimple fan extraction, strengthens structural component bounds, and isolates exact surgery losses. The unrestricted full-gain successor and unrestricted numerical upper problem remain open.

Section~\ref{sec:geometry} states the geometric definitions and the certificate-to-cell bridge. Section~\ref{sec:universal} proves the uniform refinement; Section~\ref{sec:finite} records finite cases and provenance. The following sections develop extension and family results, local upper-bound ingredients, and computational searches. The final sections describe the formal trust boundary and future obligations. The appendices contain the complete finite tables and a map from mathematical claims to proof sources. The code and evidence are available at \repo~\cite{zarzueloRepository}.
